% ===== PRÉAMBULE CANONIQUE — à coller VERBATIM en tête de CHAQUE .tex =====
\documentclass[11pt,a4paper]{article}
\usepackage[utf8]{inputenc}\usepackage[T1]{fontenc}\usepackage{lmodern}
\usepackage[french]{babel}
\usepackage{amsmath,amssymb,mathtools}\usepackage{icomma}
\usepackage{siunitx}\sisetup{locale=FR}  % \qty{}{}, \si{}, \ang{}
\usepackage{xcolor}\usepackage{colortbl}
\usepackage{tikz}\usepackage{pgfplots}\pgfplotsset{compat=1.18}
\usepackage[siunitx,europeanresistors]{circuitikz} % circuits (élec.)
\usepackage[most]{tcolorbox}\usepackage{enumitem}\usepackage{array,booktabs}
\usepackage[a4paper,margin=2.2cm]{geometry}
\usetikzlibrary{arrows.meta,calc,angles,quotes,positioning,patterns,%
  backgrounds,shapes.geometric,decorations.pathmorphing,decorations.markings,intersections}
\definecolor{primary}{RGB}{222,49,99}\definecolor{primarydark}{RGB}{150,22,55}
\definecolor{defcol}{RGB}{0,114,178}\definecolor{methcol}{RGB}{52,92,148}
\definecolor{excol}{RGB}{0,135,90}\definecolor{attncol}{RGB}{200,40,55}
\tcbset{boxbase/.style={enhanced,breakable,boxrule=0.9pt,arc=2.5pt,
  left=3mm,right=3mm,top=2.3mm,bottom=2.3mm,fonttitle=\bfseries,coltitle=white}}
% --- boîtes TUTORIEL (noms EXACTS, ne pas renommer) ---
\newtcolorbox{cle}[1][]{boxbase,colback=defcol!6,colframe=defcol,
  title={\textsc{À retenir}\ifx\relax#1\relax\relax\else\textmd{ : \itshape #1}\fi}}
\newtcolorbox{methode}[1][]{boxbase,colback=methcol!7,colframe=methcol,sharp corners,
  title={\textsc{Méthode}\ifx\relax#1\relax\relax\else\textmd{ --- \itshape #1}\fi}}
\newtcolorbox{exemplebox}[1][]{boxbase,colback=excol!5,colframe=excol,
  title={\textsc{Exemple}\ifx\relax#1\relax\relax\else\textmd{ : \itshape #1}\fi}}
\newtcolorbox{piege}[1][]{boxbase,colback=attncol!6,colframe=attncol,
  title={\textsc{Piège}\ifx\relax#1\relax\relax\else\textmd{ : \itshape #1}\fi}}
\newtcolorbox{remarque}[1][]{boxbase,colback=black!4,colframe=black!45,
  title={\textsc{Remarque}\ifx\relax#1\relax\relax\else\textmd{ : \itshape #1}\fi}}
% --- boîtes EXERCICE / CORRIGÉ (noms EXACTS, auto-comptées) ---
\newtcolorbox[auto counter]{exo}[1][]{boxbase,colback=primary!4,colframe=primary,
  title={\textsc{Exercice}~\thetcbcounter\ifx\relax#1\relax\relax\else\textmd{ --- \itshape #1}\fi}}
\newtcolorbox[auto counter]{corrige}[1][]{boxbase,colback=excol!4,colframe=excol,
  title={\textsc{Corrigé}~\thetcbcounter\ifx\relax#1\relax\relax\else\textmd{ --- \itshape #1}\fi}}
\newcommand{\vv}[1]{\vec{#1}}\newcommand{\ds}{\displaystyle}
\newcommand{\R}{\mathbb{R}}\newcommand{\N}{\mathbb{N}}\newcommand{\Z}{\mathbb{Z}}
% (un \usepackage{fancyhdr}+en-tête est possible ; il est ignoré côté web)
% =========================================================================
\begin{document}

\begin{center}
{\large\bfseries\color{primarydark} Thème 1 — Niveau Facile}\\[-2pt]
{\color{primary}\rule{5cm}{1pt}}
\end{center}

\begin{exo}[période et fréquence]
Deux ondes sont décrites par leurs caractéristiques temporelles.
\begin{enumerate}[label=\alph*)]
  \item Une onde a une fréquence $f=\SI{50}{\hertz}$. Calcule sa période $T$.
  \item Une autre onde a une période $T=\SI{4}{\milli\second}$. Calcule sa fréquence $f$
        (en \si{\hertz}).
\end{enumerate}
\end{exo}

\begin{exo}[lire la période sur un graphe]
On enregistre l'élongation $y$ d'un point du milieu \emph{en fonction du temps}.
\begin{center}
\begin{tikzpicture}
\begin{axis}[
  width=11cm,height=5.2cm,
  xmin=0,xmax=1.5,ymin=-3.6,ymax=3.6,
  xtick={0,0.25,0.5,0.75,1,1.25,1.5},
  ytick={-3,0,3},
  grid=both,grid style={black!12},axis lines=left,
  xlabel={$t$ (s)},ylabel={$y$ (cm)},
  xlabel style={at={(axis description cs:0.5,-0.12)}},
  tick label style={font=\footnotesize},samples=200,clip=false]
  \addplot[primary,line width=1.3pt,domain=0:1.5]{3*sin(deg(2*pi*x/0.5))};
\end{axis}
\end{tikzpicture}
\end{center}
\begin{enumerate}[label=\alph*)]
  \item Lis la période $T$ de l'onde sur le graphe.
  \item Déduis-en la fréquence $f$.
  \item Quelle est l'amplitude $A$ de l'onde ?
\end{enumerate}
\end{exo}

\begin{exo}[lire la longueur d'onde sur un graphe]
Voici cette fois une \emph{photographie} de l'onde à un instant donné : l'élongation $y$
est tracée \emph{en fonction de la position} $x$ le long du milieu.
\begin{center}
\begin{tikzpicture}
\begin{axis}[
  width=11cm,height=5.2cm,
  xmin=0,xmax=12,ymin=-2.6,ymax=2.6,
  xtick={0,2,4,6,8,10,12},
  ytick={-2,0,2},
  grid=both,grid style={black!12},axis lines=left,
  xlabel={$x$ (m)},ylabel={$y$ (cm)},
  xlabel style={at={(axis description cs:0.5,-0.12)}},
  tick label style={font=\footnotesize},samples=240,clip=false]
  \addplot[defcol,line width=1.3pt,domain=0:12]{2*sin(deg(2*pi*x/4))};
\end{axis}
\end{tikzpicture}
\end{center}
\begin{enumerate}[label=\alph*)]
  \item Lis la longueur d'onde $\lambda$ sur le graphe.
  \item Quelle est l'amplitude $A$ ?
\end{enumerate}
\end{exo}

\begin{exo}[relier $\lambda$, $v$ et $f$]
On travaille sur des ondes dans différents milieux.
\begin{enumerate}[label=\alph*)]
  \item Une onde de fréquence $f=\SI{200}{\hertz}$ se propage à $v=\SI{340}{\metre\per\second}$
        dans l'air. Calcule sa longueur d'onde $\lambda$.
  \item Une onde de longueur d'onde $\lambda=\SI{2}{\metre}$ et de fréquence $f=\SI{5}{\hertz}$
        avance sur l'eau. Calcule sa célérité $v$.
\end{enumerate}
\end{exo}

\begin{exo}[longitudinale ou transversale ?]
Deux dispositifs produisent une onde. En haut, on comprime périodiquement un long ressort~;
en bas, on secoue verticalement l'extrémité d'une corde tendue horizontalement.
\begin{center}
\begin{tikzpicture}[scale=1.0]
  % --- ressort (onde longitudinale) ---
  \begin{scope}[yshift=1.4cm]
    \node[left] at (-0.2,0) {(1)};
    \draw[black!75,line width=0.7pt,decorate,
          decoration={coil,aspect=0.6,segment length=2.2pt,amplitude=4pt}](0,0)--(1.6,0);
    \draw[black!75,line width=0.7pt,decorate,
          decoration={coil,aspect=0.6,segment length=5pt,amplitude=4pt}](1.6,0)--(8,0);
    \draw[->,>=Stealth,primary,line width=1pt](0.3,0.5)--(1.3,0.5)
          node[midway,above,black]{$\vv v$};
  \end{scope}
  % --- corde (onde transversale) ---
  \begin{scope}[yshift=-0.6cm]
    \node[left] at (-0.2,0) {(2)};
    \draw[defcol,line width=1.6pt,smooth,samples=120,domain=0:8]
          plot(\x,{0.45*sin(deg(2*pi*\x/2.6))});
    \draw[->,>=Stealth,black!70](-0.05,-0.75)--(-0.05,0.75)node[left]{$y$};
    \draw[->,>=Stealth,primary,line width=1pt](3.2,0.95)--(4.2,0.95)
          node[midway,above,black]{$\vv v$};
  \end{scope}
\end{tikzpicture}
\end{center}
Pour chacun des deux cas, indique s'il s'agit d'une onde \textbf{longitudinale} ou
\textbf{transversale}, et justifie en une phrase.
\end{exo}

\end{document}
